\chapter{Le monde quantique~: photons et niveaux d'énergie}\label{ch:g12:quantum-world}

Une lampe à chaleur \omterm{def:g10:light-spectra:wavelength}{infrarouge} peut baigner votre peau
toute la soirée sans la marquer~; dix minutes de soleil de montagne
\omterm{def:g11:lenses-and-eye:lens}{fin} laissent une brûlure. La lampe livre bien plus
d'\omterm{def:g10:energy-conservation:energy}{énergie} --- mais l'\omterm{def:g10:energy-conservation:energy}{énergie} n'est
pas ce qui marque la peau. La lumière, montre ce chapitre, arrive en
grains, et seul un grain qui frappe assez fort fait de la chimie~: un
grain \omterm{def:g10:light-spectra:wavelength}{ultraviolet} casse une liaison qu'un milliard de
grains \omterm{def:g10:light-spectra:wavelength}{infrarouges} ne font que réchauffer. Suivez les
grains et un siècle s'ouvre~: des métaux qui fuient des électrons, des
\omterm{def:g11:fundamental-interactions:ladder}{atomes} qui chantent en
\omterm{def:g10:light-spectra:emission}{raies spectrales}, des électrons qui interfèrent comme
des \omterm{def:g10:signals-and-waves:wave}{ondes}.

\section{La lumière arrive en grains~: le photon}

\begin{definition}[Le photon]\label{def:g12:quantum-world:photon}
La lumière de \omterm{def:g12:oscillators-and-time:oscillator}{fréquence} $\nu$ n'échange de
l'\omterm{def:g10:energy-conservation:energy}{énergie} avec la matière que par grains entiers
appelés \emph{photons}\index{photon}, chacun d'\omterm{def:g10:energy-conservation:energy}{énergie}
\[
E = h\nu = \frac{hc}{\lambda}, \qquad h = \qty{6.63e-34}{J.s},
\]
où $h$ est la \emph{constante de Planck}\index{constante de Planck}
(Planck, 1900~; Einstein, 1905) et $\lambda$ la
\omterm{def:g12:light-as-wave:lightwave}{longueur d'onde dans le vide}. L'intensité d'un
faisceau mesure le \emph{nombre} de photons par seconde~;
l'\omterm{def:g10:energy-conservation:energy}{énergie} de chaque grain est fixée par la
\omterm{def:g12:oscillators-and-time:oscillator}{fréquence} seule.
\end{definition}

\begin{notation}[L'électronvolt]\label{not:g12:quantum-world:ev}
La physique atomique utilise l'\emph{électronvolt}\index{électronvolt},
$\qty{1}{eV} = \qty{1.60e-19}{J}$~: l'\omterm{def:g10:energy-conservation:energy}{énergie} qu'une
\omterm{def:g11:fundamental-interactions:elementary}{charge élémentaire} gagne en traversant
\qty{1}{V} (\cref{ch:g11:electric-gravitational-fields}). Pratique~:
$E \,(\unit{eV}) \approx 1240/\lambda\,(\unit{nm})$.
\end{notation}

\begin{omfigure}
\begin{tikzpicture}[font=\small, x=1.48cm]
  \draw[-Stealth] (7.55,0) -- (-0.55,0);
  \draw[-Stealth, omDef] (4.9,-0.85) -- (2.3,-0.85)
    node[midway, below] {l'énergie par photon croît};
  \foreach \x/\l/\e/\n in {0/{\qty{1}{pm}}/{\qty{1.2}{MeV}}/{$\gamma$},
      1.2/{\qty{0.1}{nm}}/{\qty{12}{keV}}/{X},
      2.4/{\qty{100}{nm}}/{\qty{12}{eV}}/{UV},
      3.6/{\qty{550}{nm}}/{\qty{2.3}{eV}}/{visible},
      4.8/{\qty{10}{\micro m}}/{\qty{0.12}{eV}}/{IR},
      6.0/{\qty{1}{cm}}/{\qty{0.12}{meV}}/{micro},
      7.2/{\qty{3}{m}}/{\qty{0.4}{\micro eV}}/{radio}}{
    \draw (\x,0.07) -- (\x,-0.07);
    \node[above, omExo] at (\x,0.42) {\n};
    \node[above, font=\scriptsize] at (\x,0.10) {\l};
    \node[below, font=\scriptsize, omDef] at (\x,-0.10) {\e};}
\end{tikzpicture}

{\small Une seule échelle pour la famille électromagnétique
(\cref{ch:g12:light-as-wave})~: \omterm{def:g12:mechanical-waves:wavelength}{longueur d'onde}
(haut), \omterm{def:g10:energy-conservation:energy}{énergie} par \omterm{def:g12:quantum-world:photon}{photon}
$hc/\lambda$ (bleu), des grains radio d'un microélectronvolt aux grains
gamma de \unit{MeV} du \cref{ch:g11:nucleus-radioactivity} --- pourtant un
pointeur de \qty{1.0}{mW} verse $\num{3e15}$~grains par seconde~: la
lumière du quotidien coule aussi régulièrement que du sable versé.}
\end{omfigure}

\section{L'effet photoélectrique}

Éclairez une plaque métallique propre sous vide et des électrons
sautent --- l'\emph{effet photoélectrique}\index{effet photoélectrique}.
Ses règles ont brisé la physique classique.

\begin{proposition}[Ce que montre l'expérience]\label{prop:g12:quantum-world:facts}
Pour chaque métal il existe une \emph{fréquence seuil}\index{fréquence
seuil} $\nu_0$~:
\begin{enumerate}
  \item sous $\nu_0$, aucun électron ne part, quelle que soit
        l'intensité de la lumière~;
  \item au-dessus de $\nu_0$, l'émission démarre sans
        \omterm{def:g12:mechanical-waves:delay}{retard}, si faible soit-elle~;
  \item l'intensité élève le \emph{nombre} d'électrons par seconde,
        pas leur \omterm{def:g11:mechanical-energy:kinetic}{énergie cinétique} maximale~; la
        \omterm{def:g12:oscillators-and-time:oscillator}{fréquence} élève $E_{k,\max}$.
\end{enumerate}
Rien de cela ne s'accorde à une \omterm{def:g10:signals-and-waves:wave}{onde} continue, que
toute couleur devrait laisser se charger --- sur des mois en lumière
faible (\cref{exo:g12:quantum-world:14}) --- et dont l'intensité devrait
fixer l'\omterm{def:g10:energy-conservation:energy}{énergie} des électrons~: les trois
prédictions échouent.
\end{proposition}
\admitted

\begin{theorem}[Relation photoélectrique d'Einstein]\label{thm:g12:quantum-world:einstein}
\index{relation photoélectrique}Un électron a besoin d'au moins le
\emph{travail d'extraction}\index{travail d'extraction} $W_0$ du métal
pour s'échapper de sa surface. Un \omterm{def:g12:quantum-world:photon}{photon} est absorbé
entier par un électron, donc
\[
h\nu = W_0 + E_{k,\max}, \qquad \text{d'où } \nu_0 = \frac{W_0}{h}.
\]
\end{theorem}

\begin{proof}
Comptabilité d'\omterm{def:g10:energy-conservation:energy}{énergie}~: $h\nu$ en entrée, $W_0$
dépensé à la sortie, le reste cinétique --- au plus $h\nu - W_0$,
moins pour les électrons partant plus profond. Sous $\nu_0$ un grain
ne peut pas payer le péage de sortie, quel que soit leur nombre.
\end{proof}

\begin{definition}[Potentiel d'arrêt]\label{def:g12:quantum-world:stopping}
Pour mesurer $E_{k,\max}$, une \omterm{def:g11:circuits-and-power:voltage}{tension} inverse freine
les électrons (\cref{ch:g11:electric-gravitational-fields})~; le
\emph{potentiel d'arrêt}\index{potentiel d'arrêt} $U_s$ est la plus
petite \omterm{def:g11:circuits-and-power:voltage}{tension} qui annule le courant~: les
électrons les plus rapides échouent juste à traverser, donc
$E_{k,\max} = e\,U_s$.
\end{definition}

\begin{omfigure}
\begin{tikzpicture}
\begin{axis}[omaxis, width=8.8cm, height=5.4cm,
    xlabel={$\nu$ ($10^{14}\,\unit{Hz}$)}, ylabel={$E_{k,\max}$ (\unit{eV})},
    xmin=0, xmax=13.4, ymin=-3.1, ymax=3.2,
    xtick={2,4,6,8,10,12}, ytick={-2,-1,1,2,3}]
  \addplot[omExo, dashed, domain=0:5.55, samples=2] {0.414*x-2.3};
  \addplot[omThm, thick, domain=5.55:12.4, samples=2] {0.414*x-2.3};
  \addplot[only marks, mark=*, mark size=1.5pt, omDef] coordinates
    {(7,0.60) (8,1.01) (9,1.43) (10,1.84) (11,2.26)};
  \node[below, font=\small] at (axis cs:5.55,-0.08) {$\nu_0$};
  \node[right, font=\small, omExo] at (axis cs:0.15,-2.55) {$-W_0$};
  \node[font=\small, omThm, rotate=22] at (axis cs:10.3,1.35) {pente $= h$};
\end{axis}
\end{tikzpicture}

{\small $E_{k,\max} = eU_s$ mesuré contre la
\omterm{def:g12:oscillators-and-time:oscillator}{fréquence} (sodium)~: une droite de pente
$h$, coupant l'axe en $\nu_0$~; extrapolée (tireté), elle révèle
$-W_0$ en $\nu = 0$.}
\end{omfigure}

\begin{method}[Mesurer la constante de Planck]\label{met:g12:quantum-world:measure-h}
\begin{enumerate}
  \item Éclairer une photocellule à travers des
        \omterm{def:g11:color-light-sources:subtractive}{filtres} de
        \omterm{def:g12:mechanical-waves:wavelength}{longueurs d'onde} connues~; calculer chaque
        $\nu = c/\lambda$~; enregistrer chaque
        \omterm{def:g12:quantum-world:stopping}{potentiel d'arrêt}.
  \item Tracer $E_{k,\max} = eU_s$ en fonction de $\nu$~: Einstein
        promet la droite $h\nu - W_0$.
  \item Lire la pente~: c'\emph{est} $h$~; les intercepts donnent
        $\nu_0$ et $-W_0$.
\end{enumerate}
\end{method}

\section{Niveaux d'énergie~: pourquoi les atomes émettent des raies}

\begin{definition}[Niveaux d'énergie]\label{def:g12:quantum-world:levels}
L'\omterm{def:g10:energy-conservation:energy}{énergie} d'un électron lié dans un
\omterm{def:g11:fundamental-interactions:ladder}{atome} est \emph{quantifiée}\index{quantification}~:
restreinte à une échelle discrète de \emph{niveaux
d'énergie}\index{niveau d'énergie} $E_1 < E_2 < \dots < 0$, comptés
négatifs depuis l'électron libre. Le barreau le plus bas est
l'\emph{état fondamental}\index{état fondamental}, les autres des
\emph{états excités}\index{état excité}~; atteindre $0$ est
l'\emph{ionisation}\index{ionisation}.
\end{definition}

\begin{proposition}[L'échelle de l'hydrogène]\label{prop:g12:quantum-world:hydrogen}
Les \omterm{def:g12:quantum-world:levels}{niveaux d'énergie} de l'\omterm{def:g11:fundamental-interactions:ladder}{atome}
d'hydrogène sont
\[
E_n = -\frac{\qty{13.6}{eV}}{n^2}, \qquad n = 1, 2, 3, \dots
\]
donc $E_1 = \qty{-13.6}{eV}$, $E_2 = \qty{-3.40}{eV}$, $E_3 =
\qty{-1.51}{eV}$, $E_4 = \qty{-0.85}{eV}$, se serrant vers $0$.
\end{proposition}
\admitted

\begin{remark}[D'où vient l'échelle]\label{rem:g12:quantum-world:ladder-origin}
Pourquoi ces nombres est de la mécanique quantique, gardée pour les
volumes universitaires --- comme la formule de de Broglie~:
l'\omterm{def:g12:quantum-world:debroglie}{onde de matière} de l'électron
(\cref{def:g12:quantum-world:debroglie} ci-dessous) doit se refermer sur
elle-même autour du \omterm{def:g11:fundamental-interactions:ladder}{noyau}, et seules
certaines \omterm{def:g12:mechanical-waves:wavelength}{longueurs d'onde} tiennent --- comme une
corde de guitare, un \omterm{def:g11:fundamental-interactions:ladder}{atome} a un ensemble
discret de notes.
\end{remark}

\begin{theorem}[Photon d'une transition]\label{thm:g12:quantum-world:transition}
Un \omterm{def:g11:fundamental-interactions:ladder}{atome} ne change de niveau qu'en émettant
ou absorbant un \omterm{def:g12:quantum-world:photon}{photon} portant exactement la
différence~: descendre du niveau $p$ à un niveau plus bas $n$ émet
$h\nu = E_p - E_n$~; monter de $n$ à $p$ absorbe la même
\omterm{def:g10:energy-conservation:energy}{énergie}.
\end{theorem}

\begin{proof}
Conservation de l'\omterm{def:g10:energy-conservation:energy}{énergie}, un grain à la fois~: le
\omterm{def:g12:quantum-world:photon}{photon} est la seule monnaie que la lumière accepte, et
les livres de l'\omterm{def:g11:fundamental-interactions:ladder}{atome} doivent s'équilibrer.
\end{proof}

\begin{remark}[Spectres de raies expliqués]\label{rem:g12:quantum-world:lines}
Voici la réponse promise par le \cref{ch:g10:light-spectra}~: un gaz
chaud n'émet que les fréquences $(E_p - E_n)/h$ de sa propre échelle
--- un code-barres de raies brillantes --- et un gaz froid vole
exactement ces fréquences à la lumière \omterm{def:g10:light-spectra:white}{blanche},
imprimant des raies d'absorption sombres correspondantes. Chaque
élément a sa propre échelle, donc son propre code-barres --- comment
on lit la lumière des étoiles.
\end{remark}

\begin{example}[Les raies visibles de l'hydrogène]\label{ex:g12:quantum-world:balmer}
Les transitions aboutissant à $n = 2$ tombent dans le visible. Avec
$E \,(\unit{eV}) = 1240/\lambda\,(\unit{nm})$~: $3 \to 2$~:
\qty{1.89}{eV}, \qty{656}{nm} (rouge)~; $4 \to 2$~: \qty{2.55}{eV},
\qty{486}{nm} (bleu-vert)~; $5 \to 2$~: \qty{2.86}{eV}, \qty{434}{nm}
(violet). Les transitions vers $n = 1$ dépassent \qty{10}{eV}~: tout
\omterm{def:g10:light-spectra:wavelength}{ultraviolet}.
\end{example}

\begin{omfigure}
\begin{tikzpicture}[font=\small, y=1cm]
  \foreach \y/\n/\e in {0/1/-13.6, 4.50/2/-3.40, 5.33/3/-1.51}{
    \draw[thick] (0,\y) -- (4.2,\y);
    \node[left] at (0,\y) {$n=\n$};
    \node[right, font=\scriptsize] at (4.2,\y) {\qty{\e}{eV}};}
  \draw[thick] (0,5.63) -- (4.2,5.63) node[left,at start,yshift=-2.5pt] {$n=4$};
  \draw[thick] (0,5.76) -- (4.2,5.76) node[left,at start,yshift=3.5pt] {$n=5$};
  \draw[dashed] (0,6.0) -- (4.2,6.0) node[left, at start] {$\infty$}
    node[right, font=\scriptsize] {\qty{0}{eV} (ionisé)};
  \foreach \x/\ys/\col/\l in {1.0/5.33/omThm/656, 1.9/5.63/omDef/486,
      2.8/5.76/violet/434} \draw[-Stealth, \col, thick] (\x,\ys) --
    (\x,4.50) node[midway, left, font=\scriptsize] {\qty{\l}{nm}};
  \draw[-Stealth, omExo, thick] (3.7,4.50) -- (3.7,0)
    node[midway, right, font=\scriptsize] {\qty{122}{nm} (UV)};
\end{tikzpicture}

{\small L'échelle de l'hydrogène, à l'échelle~: les raies visibles
aboutissent toutes à $n = 2$~; la chute vers l'\omterm{def:g12:quantum-world:levels}{état
fondamental} est de l'\omterm{def:g10:light-spectra:wavelength}{ultraviolet} profond.}
\end{omfigure}

\begin{remark}[Les lasers, en un souffle]\label{rem:g12:quantum-world:laser}
Un \omterm{def:g12:quantum-world:photon}{photon} d'exactement $E_p - E_n$ passant un
\omterm{def:g11:fundamental-interactions:ladder}{atome} \emph{excité} peut le déclencher à
émettre un \omterm{def:g12:quantum-world:photon}{photon} identique --- même
\omterm{def:g10:energy-conservation:energy}{énergie}, même direction, en phase. Gardez plus
d'\omterm{def:g11:fundamental-interactions:ladder}{atomes} en haut qu'en bas et un
\omterm{def:g12:quantum-world:photon}{photon} devient une avalanche de clones~:
\emph{émission stimulée}\index{émission stimulée} --- et un
\omterm{def:g11:color-light-sources:lamps}{laser} n'est guère plus.
\end{remark}

\section{Ondes de matière}

\begin{proposition}[Les électrons interfèrent]\label{prop:g12:quantum-world:interference}
Envoyez des électrons un par un à travers une double fente
(\cref{ch:g12:light-as-wave})~: chacun atterrit comme un point net
unique --- une particule. Mais à mesure que des milliers
s'accumulent, les points dessinent des \omterm{def:g12:light-as-wave:fringes}{franges
d'interférence}~: chaque électron, seul, passe comme une
\omterm{def:g10:signals-and-waves:wave}{onde} à travers \emph{les deux} fentes. Les cristaux
diffractent aussi les faisceaux d'électrons (Davisson et Germer,
1927).
\end{proposition}
\admitted

\begin{omfigure}
\begin{tikzpicture}[font=\small]
  \foreach \o in {0,4.0,8.0} \draw[omExo] (\o,0) rectangle +(3.4,2.2);
  \foreach \x/\y in {0.62/1.71, 1.68/0.42, 2.31/1.95, 1.05/1.22,
      2.98/0.66, 1.74/1.58, 0.38/0.35, 2.40/0.98, 1.62/1.02, 0.97/1.86,
      2.95/1.60, 1.71/0.20, 1.12/0.61, 2.33/0.31}
    \fill (\x,\y) circle (0.8pt);
  \foreach \c/\f in {0.40/610, 1.05/410, 1.70/530, 2.35/350, 3.00/470}
    \foreach \y in {0.25,0.62,...,2.00}
      \fill ({4.0+\c+0.14*sin(\f*\y)},\y) circle (0.8pt);
  \foreach \x/\y in {4.72/1.42, 5.38/0.52, 6.06/1.12, 6.66/1.72,
      5.44/1.88, 7.28/0.78} \fill (\x,\y) circle (0.8pt);
  \foreach \c/\f in {0.40/610, 1.05/410, 1.70/530, 2.35/350, 3.00/470}
    \foreach \y in {0.12,0.27,...,2.10}
      \fill ({8.0+\c+0.11*sin(\f*\y)},\y) circle (0.8pt);
  \foreach \c/\f in {0.40/260, 1.05/340, 1.70/205, 2.35/290, 3.00/385}
    \foreach \y in {0.20,0.50,...,2.05}
      \fill ({8.0+\c+0.06*sin(\f*\y+40)},\y) circle (0.8pt);
  \node[below] at (1.7,-0.05) {$\sim 30$ électrons};
  \node[below] at (5.7,-0.05) {$\sim 200$}; \node[below] at (9.7,-0.05) {$\sim 2000$};
\end{tikzpicture}

{\small Un électron, un point~; des milliers, des franges. Où atterrit
le prochain point est aléatoire --- comme la désintégration d'un
\omterm{def:g11:fundamental-interactions:ladder}{noyau} (\cref{ch:g11:nucleus-radioactivity})
--- mais l'\omterm{def:g10:signals-and-waves:wave}{onde} fixe la \emph{probabilité}~: frange
claire, probable~; sombre, presque jamais. Le hasard individuel régi
par une \omterm{def:g10:signals-and-waves:wave}{onde} exacte~: le cœur de la physique
quantique.}
\end{omfigure}

\begin{definition}[Longueur d'onde de de Broglie]\label{def:g12:quantum-world:debroglie}
Toute particule de \omterm{def:g12:newtons-laws:momentum}{quantité de mouvement} $p = mv$
voyage comme une \emph{onde de matière}\index{onde de matière} de
\emph{longueur d'onde de de Broglie}\index{longueur d'onde de de
Broglie}
\[
\lambda = \frac{h}{p} = \frac{h}{mv}.
\]
\end{definition}
\admitted

\begin{example}[Pourquoi vous ne diffractez jamais]\label{ex:g12:quantum-world:never}
Un électron accéléré à travers \qty{100}{V} a $E_k = \qty{1.60e-17}{J}$,
$p = \sqrt{2m_e E_k} = \qty{5.4e-24}{kg.m/s}$, donc $\lambda \approx
\qty{0.12}{nm}$ --- l'espacement des
\omterm{def:g11:fundamental-interactions:ladder}{atomes} dans un cristal, qui agit donc comme
sa double fente. Une balle de tennis de \qty{58}{g} à \qty{25}{m/s}~:
$\lambda = 6.63\times10^{-34}/1.45 \approx \qty{4.6e-34}{m}$, quelque
$10^{19}$~fois plus petit qu'un
\omterm{def:g11:fundamental-interactions:ladder}{noyau} --- aucune fente ne la montera jamais
en franges. $h$ est si petit que le grain quantique n'affleure qu'à
l'échelle de l'\omterm{def:g11:fundamental-interactions:ladder}{atome}.
\end{example}

\begin{example}[Le microscope électronique]\label{ex:g12:quantum-world:microscope}
La \omterm{def:g12:light-as-wave:diffraction}{diffraction} interdit de résoudre un détail
beaucoup plus petit que la \omterm{def:g12:mechanical-waves:wavelength}{longueur d'onde}
utilisée (\cref{ch:g12:light-as-wave})~: la lumière visible s'arrête
près de \qty{0.5}{\micro m}. Les électrons à \qty{10}{kV} ont
$\lambda \approx \qty{12}{pm}$ --- des dizaines de milliers de fois
plus courte~: les microscopes électroniques imagent virus, machinerie
cellulaire, même des \omterm{def:g11:fundamental-interactions:ladder}{atomes} seuls.
L'\omterm{def:g12:quantum-world:debroglie}{onde de matière} a tendu à la physique un nouvel
œil.
\end{example}

\section{Exercices}

\begin{exercise}[$\star$]\label{exo:g12:quantum-world:1}
Un pointeur \omterm{def:g11:color-light-sources:lamps}{laser} vert émet \qty{1.0}{mW} à
\qty{530}{nm}. Calculer l'\omterm{def:g10:energy-conservation:energy}{énergie} d'un
\omterm{def:g12:quantum-world:photon}{photon} en \omterm{def:g11:work-of-force:work}{joules} et
électronvolts, puis les \omterm{def:g12:quantum-world:photon}{photons} par seconde.
\end{exercise}

\begin{exercise}[$\star$]\label{exo:g12:quantum-world:2}
Une station FM émet à \qty{100}{MHz}. Calculer l'\omterm{def:g10:energy-conservation:energy}{énergie}
de son \omterm{def:g12:quantum-world:photon}{photon} en \unit{eV} et la comparer à un
\omterm{def:g12:quantum-world:photon}{photon} visible~: pourquoi le grain radio est-il
indétectable~?
\end{exercise}

\begin{exercise}[$\star$]\label{exo:g12:quantum-world:3}
Le césium a $W_0 = \qty{1.9}{eV}$. Calculer sa
\omterm{prop:g12:quantum-world:facts}{fréquence seuil} et sa
\omterm{def:g12:mechanical-waves:wavelength}{longueur d'onde}. Un
\omterm{def:g11:color-light-sources:lamps}{laser} rouge à \qty{633}{nm} éjecte-t-il des
électrons du césium --- et si oui, avec quelle
\omterm{def:g11:mechanical-energy:kinetic}{énergie cinétique} maximale~?
\end{exercise}

\begin{exercise}[$\star$]\label{exo:g12:quantum-world:4}
En utilisant $E_n = -13.6/n^2$ \unit{eV}, calculer
l'\omterm{def:g10:energy-conservation:energy}{énergie} et la
\omterm{def:g12:mechanical-waves:wavelength}{longueur d'onde} du
\omterm{def:g12:quantum-world:photon}{photon} émis dans la transition $3 \to 2$ de
l'hydrogène. Sa couleur~?
\end{exercise}

\begin{exercise}[$\star$]\label{exo:g12:quantum-world:5}
À partir du graphe $E_{k,\max}$--$\nu$ du cours, lire la
\omterm{prop:g12:quantum-world:facts}{fréquence seuil} et le
\omterm{thm:g12:quantum-world:einstein}{travail d'extraction} du sodium (en
\unit{eV})~; quelle constante est la pente~?
\end{exercise}

\begin{exercise}[$\star\star$]\label{exo:g12:quantum-world:6}
De la lumière de \omterm{def:g12:mechanical-waves:wavelength}{longueur d'onde} \qty{400}{nm}
tombe sur le potassium ($W_0 = \qty{2.3}{eV}$). Calculer
l'\omterm{def:g10:energy-conservation:energy}{énergie} du \omterm{def:g12:quantum-world:photon}{photon},
l'\omterm{def:g11:mechanical-energy:kinetic}{énergie cinétique} maximale des électrons (en
\unit{eV} et \omterm{def:g11:work-of-force:work}{joules}), et le
\omterm{def:g12:quantum-world:stopping}{potentiel d'arrêt}.
\end{exercise}

\begin{exercise}[$\star\star$]\label{exo:g12:quantum-world:7}
Une photocellule est éclairée au-dessus du seuil. Que devient (a)~le
courant, (b)~le \omterm{def:g12:quantum-world:stopping}{potentiel d'arrêt}, lorsque
l'intensité est doublée à \omterm{def:g12:oscillators-and-time:oscillator}{fréquence} fixée~?
Lorsque la \omterm{def:g12:oscillators-and-time:oscillator}{fréquence} est élevée~? Quel fait
contredit le tableau d'\omterm{def:g10:signals-and-waves:wave}{onde}, et pourquoi~?
\end{exercise}

\begin{exercise}[$\star\star$]\label{exo:g12:quantum-world:8}
Quelle \omterm{def:g10:energy-conservation:energy}{énergie} minimale de
\omterm{def:g12:quantum-world:photon}{photon} ionise un
\omterm{def:g11:fundamental-interactions:ladder}{atome} d'hydrogène depuis son
\omterm{def:g12:quantum-world:levels}{état fondamental}~? Depuis $n = 2$~? Calculer les
deux \omterm{def:g12:mechanical-waves:wavelength}{longueurs d'onde} seuil et nommer leurs
domaines spectraux.
\end{exercise}

\begin{exercise}[$\star\star$]\label{exo:g12:quantum-world:9}
De la lumière \omterm{def:g10:light-spectra:white}{blanche} traverse un bocal d'hydrogène
froid (tous les \omterm{def:g11:fundamental-interactions:ladder}{atomes} dans l'\omterm{def:g12:quantum-world:levels}{état
fondamental}). Expliquer pourquoi le
\omterm{def:g12:sound-acoustics:timbre}{spectre} transmis ne montre \emph{aucune} raie
sombre dans le visible, alors que des raies d'absorption de Balmer
apparaissent dans les spectres d'étoiles chaudes.
\end{exercise}

\begin{exercise}[$\star\star$]\label{exo:g12:quantum-world:10}
Calculer la \omterm{def:g12:quantum-world:debroglie}{longueur d'onde de de Broglie} de
(a)~un électron accéléré à travers \qty{100}{V}, (b)~une mouche de
\qty{0.10}{g} à \qty{1.0}{m/s}. Comparer chacune à la taille d'un
\omterm{def:g11:fundamental-interactions:ladder}{atome} ($\sim \qty{e-10}{m}$)~; conclure.
\end{exercise}

\begin{exercise}[$\star\star$]\label{exo:g12:quantum-world:11}
Un œil adapté à l'obscurité répond à environ $90$
\omterm{def:g12:quantum-world:photon}{photons} de lumière à \qty{505}{nm}. Calculer cette
\omterm{def:g10:energy-conservation:energy}{énergie}, la comparer à
l'\omterm{def:g11:mechanical-energy:kinetic}{énergie cinétique} d'un moustique de
\qty{2.5}{mg} volant à \qty{0.50}{m/s}, et interpréter le rapport.
\end{exercise}

\begin{exercise}[$\star\star\star$]\label{exo:g12:quantum-world:12}
Une photocellule donne $U_s = \qty{1.19}{V}$ à \qty{405}{nm} et
$U_s = \qty{0.40}{V}$ à \qty{546}{nm}. À partir de ces deux points
mesurer $h$, puis le \omterm{thm:g12:quantum-world:einstein}{travail d'extraction} en
\unit{eV}. Quel métal de ce chapitre est la cathode~?
\end{exercise}

\begin{exercise}[$\star\star\star$]\label{exo:g12:quantum-world:13}
Une lampe à hydrogène montre des raies à \qty{486}{nm} et
\qty{434}{nm}. Convertir chacune en \omterm{def:g10:energy-conservation:energy}{énergie} de
\omterm{def:g12:quantum-world:photon}{photon} et identifier les deux transitions en
appariant les différences des niveaux $E_n = -13.6/n^2$ \unit{eV}.
\end{exercise}

\begin{exercise}[$\star\star\star$]\label{exo:g12:quantum-world:14}
Estimation classique du \omterm{def:g12:mechanical-waves:delay}{retard}~: de la lumière faible
d'intensité \qty{1.0e-6}{W/m^2} tombe sur le potassium
($W_0 = \qty{2.3}{eV}$). Si un \omterm{def:g11:fundamental-interactions:ladder}{atome}
collectait seulement l'\omterm{def:g10:energy-conservation:energy}{énergie} traversant sa
section (rayon \qty{1.0e-10}{m}), combien de temps lui faudrait-il pour
rassembler $W_0$~? Comparer au \omterm{def:g12:mechanical-waves:delay}{retard} observé (sous
\qty{e-9}{s}).
\end{exercise}

\begin{exercise}[$\star\star\star$]\label{exo:g12:quantum-world:15}
Dans un microscope électronique les électrons sont accélérés à travers
\qty{5.0}{kV}. Calculer leur \omterm{def:g12:quantum-world:debroglie}{longueur d'onde de de
Broglie}, la comparer à la lumière verte (\qty{550}{nm}), et expliquer
via la \omterm{def:g12:light-as-wave:diffraction}{diffraction}
(\cref{ch:g12:light-as-wave}) pourquoi cela achète de la résolution.
\end{exercise}

\section{Problème~: Le labo de physique du panneau solaire}

\begin{problem}[{Problème du week-end --- le labo de physique du
panneau solaire~: les $10^{21}$~grains qu'un toit boit chaque seconde,
une photocellule qui mesure la \omterm{def:g12:quantum-world:photon}{constante de Planck},
une lampe à hydrogène lue raie par raie --- un week-end, un $h$}]
\label{pb:g12:quantum-world:1}
Votre classe instrumente le panneau solaire du toit de l'école pour un
week-end. Données~: irradiance solaire à midi \qty{1.0e3}{W/m^2}~;
\omterm{def:g12:mechanical-waves:wavelength}{longueur d'onde} solaire moyenne \qty{550}{nm}~;
aire du panneau \qty{1.0}{m^2}, \omterm{def:g10:energy-conservation:efficiency}{rendement}
$20\%$~; le jour livre l'équivalent de \qty{5.0}{h} de plein soleil.

\medskip
\noindent\textbf{Partie I --- Comptabilité des
\omterm{def:g12:quantum-world:photon}{photons}.}
\begin{enumerate}
  \item Calculer l'\omterm{def:g10:energy-conservation:energy}{énergie} d'un
        \omterm{def:g12:quantum-world:photon}{photon} solaire moyen, en
        \omterm{def:g11:work-of-force:work}{joules} et électronvolts.
  \item En déduire le nombre de \omterm{def:g12:quantum-world:photon}{photons} frappant le
        panneau par seconde à midi.
  \item Calculer la \omterm{prop:g11:circuits-and-power:power}{puissance électrique} du panneau
        à midi, puis l'\omterm{def:g10:energy-conservation:energy}{énergie} produite sur la
        journée, en \unit{kW.h}.
  \item Le courant du panneau paraît parfaitement lisse. À l'aide de
        la question~2, expliquer pourquoi l'arrivée grain par grain
        est invisible.
  \item La lumière solaire porte plus d'\omterm{def:g10:energy-conservation:energy}{énergie}
        \omterm{def:g10:light-spectra:wavelength}{infrarouge} que
        \omterm{def:g10:light-spectra:wavelength}{ultraviolet}, pourtant seul
        l'\omterm{def:g10:light-spectra:wavelength}{ultraviolet} brûle la peau (liaisons
        $\sim \qtyrange{3}{4}{eV}$). Expliquer en grains.
\end{enumerate}

\medskip
\noindent\textbf{Partie II --- Mesurer $h$ avec la photocellule du
kit.}
La photocellule du kit a une cathode inconnue. Des
\omterm{def:g11:color-light-sources:subtractive}{filtres} donnent~:
\begin{center}\small
\begin{tabular}{lcccc}
$\lambda$ (\unit{nm}) & 365 & 405 & 436 & 546\\
$U_s$ (\unit{V}) & 1.50 & 1.16 & 0.94 & 0.37
\end{tabular}
\end{center}
\begin{enumerate}[resume]
  \item Que mesure le \omterm{def:g12:quantum-world:stopping}{potentiel d'arrêt}~? Relier
        $U_s$, $\nu$, $W_0$, $h$.
  \item Convertir les quatre
        \omterm{def:g12:mechanical-waves:wavelength}{longueurs d'onde} en fréquences (en
        $10^{14}\,\unit{Hz}$).
  \item Pourquoi tracer $eU_s$ en fonction de $\nu$ doit-il donner une
        droite~? Vérifier l'alignement des points.
  \item À partir des points extrêmes, calculer la pente~: votre $h$
        mesuré.
  \item En déduire le \omterm{thm:g12:quantum-world:einstein}{travail d'extraction} de la
        cathode (\unit{eV}), la \omterm{prop:g12:quantum-world:facts}{fréquence seuil} et la
        \omterm{def:g12:mechanical-waves:wavelength}{longueur d'onde} seuil.
  \item Quelle partie de la lumière solaire pourrait éjecter des
        électrons de cette cathode --- et quelle moitié du
        \omterm{def:g12:sound-acoustics:timbre}{spectre} lui est inutile~?
  \item Comparer votre $h$ à \qty{6.63e-34}{J.s}~: erreur relative en
        pour cent~?
\end{enumerate}

\medskip
\noindent\textbf{Partie III --- La lampe d'étalonnage à hydrogène.}
L'échelle de \omterm{def:g12:mechanical-waves:wavelength}{longueur d'onde} du kit est
vérifiée contre une lampe à hydrogène, $E_n = -13.6/n^2$ \unit{eV}.
\begin{enumerate}[resume]
  \item Calculer $E_2$, $E_3$, $E_4$ et $E_5$.
  \item Calculer les \omterm{def:g12:quantum-world:photon}{énergies des photons} des transitions
        $3 \to 2$, $4 \to 2$, $5 \to 2$.
  \item En déduire les trois
        \omterm{def:g12:mechanical-waves:wavelength}{longueurs d'onde} et donner la couleur
        de chaque raie.
  \item Pourquoi la lampe émet-elle des \emph{raies} plutôt qu'un
        arc-en-ciel continu~? Une phrase, nommant l'idée clé.
  \item L'hydrogène de la lampe pourrait-il absorber de la lumière à
        \qty{500}{nm}~? Justifier avec l'échelle.
\end{enumerate}

\medskip
\noindent\textbf{Partie IV --- Inspecter les cellules~:
l'\omterm{def:g10:signals-and-waves:wave}{onde} électronique.}
Des microfissures dans une cellule se placent bien sous la portée de
\qty{0.5}{\micro m} des microscopes optiques~; le microscope
électronique à balayage du labo tourne à \qty{10}{kV}.
\begin{enumerate}[resume]
  \item Calculer la vitesse des électrons~; vérifier qu'elle reste
        sous $0.25\,c$ (formule classique).
  \item Calculer leur \omterm{def:g12:newtons-laws:momentum}{quantité de mouvement} et leur
        \omterm{def:g12:quantum-world:debroglie}{longueur d'onde de de Broglie}.
  \item La résolution évolue avec la
        \omterm{def:g12:mechanical-waves:wavelength}{longueur d'onde}~: calculer le gain sur
        la lumière verte (\qty{550}{nm})~; conclure avec les deux
        nombres du week-end --- votre $h$ et ce gain.
\end{enumerate}
\end{problem}
